\section{Transfer in die medizinische Diagnostik} \label{sec:transfer}

\subsection{Transfer Learning: Methodische Übertragung zwischen Domänen}
\label{sec:transfer_learning}

Transfer Learning ermöglicht die methodische Übertragung von Modellwissen zwischen
Anwendungsdomänen. Ein auf einem großen allgemeinen Bilddatensatz vortrainiertes Modell
verfügt über hierarchische Repräsentationen grundlegender Bildstrukturen -- Kanten,
Texturen, geometrische Formen. Diese Repräsentationen lassen sich durch gezieltes
Fine-Tuning auf medizinische Bilddaten spezialisieren. Da medizinische Bildbestände
annotationsaufwändig und im Vergleich zu allgemeinen Datensätzen deutlich kleiner sind,
überbrückt dieser Ansatz den Datenmangel systematisch \cite{transferMedical2024}.

Aktuelle Forschungsergebnisse zeigen, dass eine Vortrainingsphase auf domänenverwandten
Daten die Transferleistung weiter verbessert, da die Domänendistanz zum Zieldatensatz
sinkt \cite{aiMedicalImaging2025}.

\subsection{CheXNet: Empirischer Beleg des Transfers} \label{sec:chexnet}

Das Modell \textbf{CheXNet} der Stanford University demonstriert den realisierten Transfer
exemplarisch. Ein 121-schichtiges DenseNet wurde auf ChestX-ray14 trainiert -- einem
öffentlichen Datensatz des National Institutes of Health (NIH) mit über 100.000
annotierten Röntgenaufnahmen und 14 Krankheitsklassen. Das Modell wurde anschließend in
einem direkten Benchmark gegen vier praktizierende Fachradiologen evaluiert.

CheXNet erzielte bei der Erkennung von Pneumonie einen F1-Score von 0{,}435, gegenüber
dem Radiologen-Durchschnitt von 0{,}387 -- statistisch signifikant auf einem gemeinsamen,
verblindeten Testdatensatz \cite{chexnet2017}. Das Ergebnis belegt, dass CNN-Architekturen
aus der industriellen Defekterkennung auf medizinischen Bilddaten Klassifikationsleistungen
auf dem Niveau von Fachspezialisten erzielen.

\subsection{KI als \glqq zweiter Leser\grqq{}: Kombinierter Einsatz von Mensch und System}
\label{sec:zweiter_leser}

Das praktisch wirksamste Einsatzmodell ist nicht die vollständige Automatisierung der
Befundung, sondern der Einsatz von KI als ergänzenden zweiten Leser: Das System
übernimmt die Vorselektion, grenzwertige Fälle werden dem Radiologen zur abschließenden
klinischen Entscheidung vorgelegt. Dieses Prinzip entspricht dem in industriellen
Anwendungen etablierten Modell der aufgabenspezifischen Arbeitsteilung zwischen System
und Mensch. Die quantitativen Auswirkungen dieses Ansatzes werden in
Abschnitt~\ref{sec:ki_vs_mensch} dargelegt.
